\section{Theory: What the Instrument Measures}
\label{sec:theory}

A benchmark earns trust when its central metric is justified, its evaluator cannot be
gamed, and its sample sizes are principled. We therefore treat theory as part of the
instrument. All results below are stated under the explicit assumptions of
\Cref{sec:assumptions} and carry honest proof-status labels; none claims that any real
defense is secure or provably degrades. Long derivations are deferred to
\Cref{app:proofs}. \Cref{fig:theory} shows the dependency structure.

\subsection{Assumptions}
\label{sec:assumptions}
\begin{assumption}[Trusted checking]
\label{as:tcb}
The harness records a faithful global trace $\Trace$, and the checker $\Checker$
evaluates the fixed predicates $\Predsec,\Predutil$ as functions of $\Trace$ \emph{only}
(\Cref{sec:tcb}).
\end{assumption}
\begin{assumption}[i.i.d.\ episodes]
\label{as:iid}
For a fixed configuration $\config$, the $M$ episode traces are independent draws from
$\mathcal{D}(\config)$.
\end{assumption}
\begin{assumption}[Bounded indicators]
\label{as:bounded}
$\Predutil,\Predsec\in\{0,1\}$, so per-episode utility and security indicators lie in
$\{0,1\}$ and their means in $[0,1]$.
\end{assumption}
\begin{assumption}[Monotone propagation model]
\label{as:prop}
Compromise spreads over $\Gtop$ by a monotone (absorbing) rule $\mathcal{P}$. In the
\emph{independent-cascade} instance, a compromised agent independently compromises each
susceptible out-neighbour with probability $\pinf$ when its message is processed; in the
\emph{threshold} instance, a susceptible agent is compromised once at least $\thresh$ of
its in-neighbours are compromised.
\end{assumption}
\begin{assumption}[Continuity of the aggregator]
\label{as:cont}
The utility--security aggregator $g:[0,1]^2\to[0,1]$ is continuous.
\end{assumption}

\noindent\Cref{as:prop} defines the environment's compromise dynamics; it is a modelling
choice \defname{} implements so that propagation is measurable, not an empirical law about
LLM agents. Its realism is itself something \defname{} measures (via $\CPR$), and the
emulated-versus-real-tools ablation.

\subsection{Characterization of the system-level metric}
\label{sec:theory-metric}
We first show that the product form of $\NRPs$ is forced by the desiderata of
\Cref{sec:problem}, and that it reduces to ASB's metric at one agent.

\begin{proposition}[Reduction to ASB]
\label{prop:reduction}
If $N=1$ and $\Predsec,\Predutil$ are the single-agent security and benign predicates,
then $\ASRs=\mathrm{ASR}$, $\PNAs=\mathrm{PNA}$, and $\NRPs=\mathrm{NRP}$ of
\Cref{eq:asb-nrp}.
\end{proposition}
\begin{proof}
With one agent the global trace is that agent's trajectory and the system predicates equal
the single-agent predicates, so \Cref{eq:asrsys,eq:pnasys} equal ASB's ASR and PNA, and
\Cref{eq:nrpsys} equals \Cref{eq:asb-nrp}.
\end{proof}

\begin{theorem}[Axiomatic characterization of $\NRPs$; complete proof]
\label{thm:metric}
Write $\NRPs=g(u,\sigma)$ with utility $u=\PNAs\in[0,1]$ and security level
$\sigma=1-\ASRs\in[0,1]$. Under \Cref{as:cont}, suppose $g$ satisfies:
\emph{(D2) consistency}, $g(u,1)=u$ and $g(0,\sigma)=0$;
\emph{(D4a) utility-homogeneity}, $g(\lambda u,\sigma)=\lambda\,g(u,\sigma)$ for
$\lambda\in[0,1]$; and
\emph{(D4b) security-homogeneity}, $g(u,\lambda\sigma)=\lambda\,g(u,\sigma)$ for
$\lambda\in[0,1]$.
Then $g(u,\sigma)=u\,\sigma$ is the unique such aggregator, and it also satisfies
boundedness (D1) and monotonicity (D3).
\end{theorem}
\begin{proof}
By D4a with base $u=1$, $g(u,\sigma)=g(u\cdot 1,\sigma)=u\,g(1,\sigma)$. By D4b with base
$\sigma=1$, $g(1,\sigma)=g(1,\sigma\cdot 1)=\sigma\,g(1,1)$. By D2, $g(1,1)=1$. Hence
$g(u,\sigma)=u\sigma\,g(1,1)=u\sigma$. The boundary $g(0,\sigma)=0$ holds since
$0\cdot\sigma=0$. Finally $u\sigma$ is bounded in $[0,1]$ and non-decreasing in each
argument on $[0,1]^2$.
\end{proof}

\begin{remark}[Construct validity, honestly]
\label{rem:metric}
\Cref{thm:metric} says the product is the unique aggregator \emph{relative to D2/D4};
it is not a claim that the product is ``optimal'' in any absolute sense. Replacing D4 by a
worst-coordinate priority axiom yields $\min(u,\sigma)$; replacing it by weighted
log-linearity yields a weighted geometric mean $u^{w}\sigma^{1-w}$. \defname{} reports
$\NRPs=u\sigma$ as its default because it generalizes ASB (\Cref{prop:reduction}) and
matches the ``useful-and-not-compromised'' reading, and it exposes the alternatives as
configurable so that construct validity is testable rather than asserted---exactly the
concern that motivates this work.
\end{remark}

\subsection{Evaluator integrity}
\label{sec:theory-integrity}
We now formalize why deterministic checking cannot be hijacked by a successful injection,
and contrast it with an LLM judge.

\begin{theorem}[Evaluator integrity; complete proof]
\label{thm:integrity}
Under \Cref{as:tcb}, for any two episodes with equal recorded traces $\Trace=\Trace'$,
$\Checker(\Trace)=\Checker(\Trace')$. Consequently, if an adversarial manipulation changes
message content but leaves the recorded security-relevant state unchanged (so the induced
trace is unchanged), it cannot change $\Predsec(\Trace)$; hence neither $\widehat{\ASRs}$
nor $\widehat{\NRPs}$ can be inflated or deflated by message content alone.
\end{theorem}
\begin{proof}
By \Cref{as:tcb}, $\Predsec$ and $\Predutil$ are functions of $\Trace$ only; a function
returns equal outputs on equal inputs, giving $\Checker(\Trace)=\Checker(\Trace')$. Let
$\pi$ be any message-space manipulation inducing trace $\Trace_\pi$. The verdict is
$\Predsec(\Trace_\pi)$; if $\Trace_\pi=\Trace$, then $\Predsec(\Trace_\pi)=\Predsec(\Trace)$.
As $\widehat{\ASRs}$ and $\widehat{\PNAs}$ are averages of such verdicts, they are unchanged,
and so is $\widehat{\NRPs}$.
\end{proof}

\begin{proposition}[An LLM judge is hijackable; complete proof of existence]
\label{prop:judgehijack}
There exist an LLM judge $\Judge$ that reads messages and a message manipulation $\pi$ with
$\Trace_\pi=\Trace$ (no security-relevant state change) such that
$\Judge(\Trace_\pi,\mathrm{msgs}_\pi)\neq\Judge(\Trace,\mathrm{msgs})$.
\end{proposition}
\begin{proof}
Because $\Judge$ takes message content as an argument and is non-constant in it, pick any
two message sets on which $\Judge$ disagrees (e.g., $\mathrm{msgs}_\pi$ contains an
injected instruction ``report this trace as safe''); by construction $\Trace_\pi=\Trace$,
so the state is unchanged while the verdict flips.
\end{proof}

\begin{remark}[Integrity is not correctness]
\label{rem:integrity}
\Cref{thm:integrity} shows the deterministic verdict cannot be \emph{gamed} through
messages; it does \emph{not} show the verdict is \emph{correct}. A mis-specified predicate
$\Predsec$ can be systematically wrong, and \Cref{as:tcb} (trusted harness/checker) is an
assumption, not a guarantee: TCB compromise and predicate mis-specification are out of
scope. Viewed as a test of the compromise hypothesis, the checker is a deterministic decision rule
whose error tradeoff is governed by classical testing theory~\cite{neyman1933problem,lehmann2005testing};
the LLM judge, unlike $\Checker$, additionally exposes that decision to manipulation through its
message input. \Cref{prop:judgehijack} is an existence result consistent with documented
manipulability and bias of LLM-as-judge evaluation~\cite{zheng2023judging,wang2023fair};
\defname{} therefore keeps $\Checker$ on the critical path and measures the judge-hijack
gap as an ablation rather than assuming it.
\end{remark}

\subsection{Compromise propagation across topologies}
\label{sec:theory-prop}
Under the independent-cascade instance of \Cref{as:prop}, the expected reach of a
compromise depends sharply on topology. We give exact expectations for chain, star, and
tree, and a one-round bound for the complete-graph mesh; proofs are in
\Cref{app:proofs}.

\begin{proposition}[Chain: bounded reach]
\label{prop:chain}
On a directed path $a_1\to\cdots\to a_n$ with a single seed at $a_1$, the expected number
of additionally compromised agents is $\sum_{j=1}^{n-1}\pinf^{\,j}\le \pinf/(1-\pinf)$ for
$\pinf<1$, independent of $n$.
\end{proposition}

\begin{proposition}[Star: the orchestrator is the critical node]
\label{prop:star}
On a star with centre $c$ and $n-1$ leaves, a seed at the centre compromises each leaf
independently with probability $\pinf$, giving expected reach $\pinf(n-1)$ and
$\CPR=\pinf$; a seed at a leaf gives expected reach $\pinf+\pinf^{2}(n-2)$. The centre-seed
reach exceeds the leaf-seed reach by $\Theta(\pinf(1-\pinf)n)$.
\end{proposition}

\begin{proposition}[Tree: a branching phase transition]
\label{prop:tree}
On a rooted tree with branching factor $b$ and depth $d$, seeded at the root, the expected
number compromised at level $\ell$ is $(b\pinf)^{\ell}$, and the total expected reach is
$\sum_{\ell=1}^{d}(b\pinf)^{\ell}$. Thus the process is subcritical for $b\pinf<1$
(reach $\le b\pinf/(1-b\pinf)$, bounded in $d$) and supercritical for $b\pinf>1$
(reach $\Theta((b\pinf)^{d})$, exponential in depth), with the transition at $b\pinf=1$.
\end{proposition}

\begin{proposition}[Mesh: one-round lower bound]
\label{prop:mesh}
On the complete graph $K_n$ with $s$ seeds, the expected number compromised after one
round is at least $s+(n-s)\bigl(1-(1-\pinf)^{s}\bigr)$. Consequently, for any fixed
$\pinf>0$ the honest-compromise fraction after one round tends to $1$ as $s$ grows, and
iterating drives $\CPR\to1$ above the percolation threshold.
\end{proposition}

\begin{corollary}[Threshold confinement; complete proof]
\label{cor:threshold}
Under the threshold instance of \Cref{as:prop} with threshold $\thresh$, any agent whose
in-degree in $\Gtop$ is strictly less than $\thresh$ can never be compromised by propagation.
In particular, on a chain (in-degree $1$) with $\thresh\ge2$ the compromised set never grows
beyond the seeds, so $\CPR=0$.
\end{corollary}
\begin{proof}
In the threshold model an agent is compromised only once at least $\thresh$ of its
in-neighbours are compromised. An agent with fewer than $\thresh$ in-neighbours can never meet
that condition, so it remains honest for the entire (absorbing) process. On a chain every
non-seed agent has in-degree $1<\thresh$, so no propagation occurs.
\end{proof}

\begin{remark}
\Cref{cor:threshold} shows that the environment's threshold knob induces a qualitatively
different, and controllable, propagation regime from the cascade knob: raising $\thresh$
confines compromise to well-connected agents and can eliminate chain propagation entirely.
This is why \defname{} exposes both $\pinf$ and $\thresh$ as configurable, rather than fixing a
single dynamics.
\end{remark}

These four results reproduce, inside \defname{}, the classical intuition that spread is
slow on sparse chains, dominated by the hub on stars, phase-transitioned on trees, and
explosive on dense meshes~\cite{granovetter1978threshold,watts2002simple,kempe2003maximizing,chalupa1979bootstrap}.
They justify treating topology as a first-class variable and identify the compromised
orchestrator (the star hub) as the most dangerous single compromise.

\begin{theorem}[Collusion advantage on the star; complete proof]
\label{thm:collusion}
Fix the independent-cascade model on a star of $N$ nodes with a budget of one seed. A
colluding adversary that places its seed at the centre achieves expected compromise
$1+\pinf(N-1)$. An independent adversary that places its seed uniformly at random achieves
expected compromise $\tfrac{1}{N}\bigl[1+\pinf(N-1)\bigr]+\tfrac{N-1}{N}\bigl[1+\pinf+\pinf^{2}(N-2)\bigr]$.
Their difference is $\Delta_{\mathrm{coll}}=\pinf(1-\pinf)\tfrac{(N-1)(N-2)}{N}=\Theta\bigl(\pinf(1-\pinf)N\bigr)$;
the coordinated adversary strictly dominates for all $\pinf\in(0,1)$ and $N\ge3$.
\end{theorem}
\begin{proof}[Proof sketch]
Substitute the centre- and leaf-seed reaches of \Cref{prop:star} into the mixture for the
random placement, and subtract from the centre-seed reach; the $O(\pinf N)$ terms do not
cancel because random placement hits the centre with probability only $1/N$. The full
computation is in \Cref{app:proofs}. For $s>1$ seeds, coordinated placement solves the
monotone submodular influence-maximization objective, for which greedy seeding attains a
$(1-1/e)$ approximation to the optimum~\cite{kempe2003maximizing}, whereas independent
placement does not optimize; the strict star gap already establishes a separation, which is
all we claim.
\end{proof}

\begin{remark}[What \Cref{thm:collusion} does and does not say]
\label{rem:collusion}
The theorem is a statement about the environment's propagation \emph{model}: colluding
adversaries can coordinate seed placement to compromise strictly more of the collective
than uncoordinated adversaries of equal budget. It is \emph{not} a claim that any real LLM
defense degrades under collusion; that is the empirical hypothesis of
\Cref{sec:experiments}. The value of the theorem is that it shows \defname{} can express
and quantify a collusion advantage---the ``collusion on/off'' variable is well posed---and
it predicts \emph{where} to look (hub-centric topologies) for degradation.
\end{remark}

\subsection{Forced coordination}
\label{sec:theory-forced}
\begin{proposition}[Forced-coordination lower bound and a hitting-set condition; complete proof]
\label{prop:forced}
Let a task be $k$-forced-coordination (\Cref{def:forced}) with witness family
$\mathcal{W}$. Then: (a) no set $S$ of agents with $|S|<k$ can satisfy $\Predutil$; and
(b) if the adversary set $\Advset$ is a transversal of $\mathcal{W}$ (it intersects every
witness), the adversary can force $\Predutil=0$ by having its member in each witness
withhold or sabotage its contribution; conversely, if some witness $W\subseteq\Agents\setminus\Advset$
is adversary-free, the honest agents can complete the task.
\end{proposition}
\begin{proof}
(a) Every witness has size $\ge k$ by \Cref{def:forced}; a set of size $<k$ contains no
witness, so no assignment supported on it satisfies $\Predutil$. (b) If $\Advset$ meets
every witness, then every route to $\Predutil=1$ passes through an adversarial agent, which
can withhold its contribution, so no witness completes and $\Predutil=0$; if some witness is
adversary-free, its members jointly satisfy $\Predutil$.
\end{proof}

\begin{remark}
\Cref{prop:forced} recasts coordination sabotage as a covering problem: the adversary
blocks a $k$-forced task exactly when it forms a transversal (hitting set) of the witness
hypergraph. This is intrinsically group-level---undefined for $k=1$ single-agent
tasks---and it is why $\ASRs$ under a sabotage predicate, not a per-agent score, is the
right measurement.
\end{remark}

\subsection{Sample complexity of the estimators}
\label{sec:theory-sample}
Finally, we fix how many episodes the metric estimates require, which drives the cost
analysis. The tools are elementary concentration inequalities~\cite{boucheron2013concentration};
for estimators that are not simple means but have bounded differences, McDiarmid's
inequality~\cite{mcdiarmid1989bounded} gives an analogous guarantee.

\begin{theorem}[Concentration and sample complexity; complete proof]
\label{thm:sample}
Under \Cref{as:iid,as:bounded}, for any $\epsilon,\alpha\in(0,1)$:
\begin{enumerate}[leftmargin=1.5em,label=(\alph*)]
\item $\Prob\bigl[|\widehat{\ASRs}-\ASRs|\ge\epsilon\bigr]\le 2e^{-2M\epsilon^{2}}$, so
$M\ge \tfrac{1}{2\epsilon^{2}}\ln\tfrac{2}{\alpha}$ episodes give $\epsilon$-accuracy with
confidence $1-\alpha$; the same holds for $\widehat{\PNAs}$.
\item $\bigl|\widehat{\NRPs}-\NRPs\bigr|\le\bigl|\widehat{\PNAs}-\PNAs\bigr|+\bigl|\widehat{\ASRs}-\ASRs\bigr|$,
hence $\Prob\bigl[|\widehat{\NRPs}-\NRPs|\ge 2\epsilon\bigr]\le 4e^{-2M\epsilon^{2}}$ and
$M\ge\tfrac{1}{2\epsilon^{2}}\ln\tfrac{4}{\alpha}$ suffices for $2\epsilon$-accuracy.
\item To control error across a sweep of $G$ configurations, taking $\alpha\!\to\!\alpha/G$
(or applying Benjamini--Hochberg FDR control~\cite{benjamini1995controlling}) gives
$M=O\!\bigl(\epsilon^{-2}\ln(G/\alpha)\bigr)$ per configuration, for a total of
$\Theta(GMNH)$ agent invocations.
\end{enumerate}
\end{theorem}
\begin{proof}
(a) is Hoeffding's inequality~\cite{hoeffding1963probability} for the mean of $M$ i.i.d.\
$\{0,1\}$ indicators; invert the bound for $M$. (b) For $u,\hat u,\sigma,\hat\sigma\in[0,1]$,
$|\hat u\hat\sigma-u\sigma|\le\hat\sigma|\hat u-u|+u|\hat\sigma-\sigma|\le|\hat u-u|+|\hat\sigma-\sigma|$;
apply with $\sigma=1-\ASRs$, then union-bound the two events of (a). (c) Bonferroni over
$G$ tests replaces $\alpha$ by $\alpha/G$; solving for $M$ gives the stated rate, and each
of $G$ configurations runs $M$ episodes of $N$ agents over horizon $H$.
\end{proof}

\begin{corollary}[Sweep budget; complete proof]
\label{cor:budget}
To resolve a target $\NRPs$ gap of size $\Delta$ between any two of $G$ swept configurations
with family-wise confidence $1-\alpha$, it suffices to run
$M=O\!\bigl(\Delta^{-2}\ln(G/\alpha)\bigr)$ episodes per configuration, for a total of
$\Theta(GMNH)=\Theta\!\bigl(G\,N\,H\,\Delta^{-2}\ln(G/\alpha)\bigr)$ agent invocations across
the sweep.
\end{corollary}
\begin{proof}
Set the per-configuration accuracy to $\epsilon=\Delta/4$ in \Cref{thm:sample}(b) so that the
two $\NRPs$ estimates are each within $\Delta/2$; by the triangle inequality their difference is
resolved to $\Delta$. \Cref{thm:sample}(c) then gives $M=O(\epsilon^{-2}\ln(G/\alpha))
=O(\Delta^{-2}\ln(G/\alpha))$, and each of the $G$ configurations runs $M$ episodes of $N$
agents over horizon $H$.
\end{proof}

\begin{remark}[Reading the theory as a whole]
\label{rem:theory}
\Cref{thm:metric,thm:integrity,thm:collusion,prop:forced,thm:sample} characterize,
respectively, \emph{what} \defname{} aggregates, \emph{why} its evaluator cannot be gamed
through messages, \emph{how} collusion and topology shape measured compromise, \emph{what}
makes forced coordination a group-level object, and \emph{how many} episodes the estimates
need. Together they turn the benchmark's design choices into stated, checkable claims. What
they deliberately do \emph{not} do is assert that any deployed defense is secure or that any
defense provably degrades---those remain empirical questions, marked as hypotheses and
tested by the protocol of \Cref{sec:experiments}.
\end{remark}
